\documentclass[a4paper]{article}     
\usepackage{amsfonts}
\usepackage{amsmath}
\usepackage{amssymb}
\usepackage{graphicx}
\usepackage{cancel}

\newcommand{\asa}{\Leftrightarrow} %als slechts als
\newcommand{\adR}{\Rightarrow} %als dan Rechts
\newcommand{\adL}{\Leftarrow}
\newcommand{\Lh}{\left(} %Links haakje
\newcommand{\Rh}{\right)}
\newcommand{\Lv}{\left[} %Links vierkanthaakje
\newcommand{\Rv}{\right]}
\newcommand{\La}{\left\{} %Links acolade
\newcommand{\Ra}{\right\}}
\newcommand{\Ras}{\right.} %Rechts acolade stelsel (omdat om stelsels te maken heb je een punt nodig)
\newcommand{\Ls}{\left|} %Links (absolutewaarde)streep
\newcommand{\Rs}{\right|}
\newcommand{\pnR}{\rightarrow} %pijl naar Rechts
\newcommand{\limiet}{\lim\limits_}
\newcommand{\schrap}{\cancel}
\newcommand{\schrapb}{\bcancel} %backslash
\newcommand{\schrapk}{\xcancel} %kruis
\newcommand{\vervang}{\cancelto}
\newcommand{\intgrl}{\displaystyle\int} %integraal teken (bij het berekenen van primitieven)

\begin{document}

\noindent \large Auteur: Yoren Collier \\
\noindent \large Studierichting: Informatica\\
\noindent \large Oefeningenreeks 5A nr. 6.2\\

\medskip

\normalsize

$ BI \intgrl_0^1 \dfrac{dx}{3x+2} \adR OBI = \intgrl \dfrac{dx}{3x+2}$\\

Stel: $t = 3x +2 \adR dt = 3dx \asa dx = \dfrac{dt}{3}$\\

$= \intgrl \dfrac{1}{t} \dfrac{dt}{3} = \dfrac{1}{3} \ln \Lh 3x +2 \Rh +c$\\

$BI = \Lv \dfrac{1}{3} \ln \Ls 3x +2 \Rs \Rv_0^1 = \dfrac{1}{3} \ln \Ls 3 \cdot 1 +2 \Rs - \dfrac{1}{3} \ln \Ls 3 \cdot 0 +2 \Rs = \dfrac{1}{3} \Lv \ln \Ls 5 \Rs - \ln \Ls 2 \Rs \Rv = \dfrac{1}{3} \ln \Ls \dfrac{5}{2} \Rs$\\

\begin{thebibliography}{999}
%\bibitem{a} Referentie
\end{thebibliography}
\end{document} 